\section{A Distortion Model for Trailing-One Sign Embedding}
\label{sec:distortion}

The forward direction of the channel---choose candidates, embed, measure PSNR---does not tell a designer how many bits a frame can take before its quality drops below a target. This section derives the distortion of one flip in closed form, models its propagation, and inverts the result into a per-IDR bit budget. All energies are sums of squared errors (SSE) of 8-bit samples; $\PSNR=10\log_{10}(255^2 N_{\mathrm{px}}/\SSE)$ for a plane of $N_{\mathrm{px}}$ samples.

\subsection{Direct distortion of one flip}

\begin{lemma}[Single flip]
\label{lem:flip}
Flipping the sign of a trailing one at raster position $(i,j)$ of a $4\times4$ block coded with quantization parameter $\mathit{QP}$ changes the reconstructed residual of that block, before rounding and clipping, by an error of energy
\begin{equation}
E_{ij}=4\,\Qs(\mathit{QP})^2\,\kappa_m(i,j),
\label{eq:flip}
\end{equation}
\begin{equation}
\kappa_m(i,j)=\frac{v_m(i,j)^2\,\|b_i\|^2\|b_j\|^2}{16\,v_{m,0}^2},
\label{eq:kappa}
\end{equation}
with $m=\mathit{QP}\bmod6$. $\kappa_m(i,j)=1$ for the even/even class (including the DC position) and $\kappa_m(i,j)\in[0.925,1.057]$ for all positions and all $m$.
\end{lemma}
\begin{proof}
The level changes by $\Delta c=\mp2$. By Eq.~\eqref{eq:dequant} the scaled coefficient changes by $\Delta d=\mp2\,v_m(i,j)\,2^{\lfloor\mathit{QP}/6\rfloor}$ and no other coefficient changes. By linearity of Eq.~\eqref{eq:inverse} before the final rounding, the residual error is $\Delta e(x,y)=\Delta d\, b_i(x)\,b_j(y)/64$, whose energy is $(\Delta d/64)^2\|b_i\|^2\|b_j\|^2$. Substituting $\Qs=v_{m,0}2^{\lfloor\mathit{QP}/6\rfloor}/16$ gives Eq.~\eqref{eq:flip}. The bounds follow by evaluating the six tables: $\kappa\in\{1,1.056\}$, $\{1,1.012,1.046\}$, $\{0.925,0.947,1\}$, $\{1,1.033,1.054\}$, $\{0.954,0.977,1\}$ and $\{1,1.014,1.020\}$ for $m=0,\dots,5$.
\end{proof}

The near-constancy of $\kappa$ is not a coincidence: the scaling factors $v_m$ were designed to compensate the unequal norms of the integer basis~\cite{malvar2003transform}, so the scaled transform is close to orthonormal and, by Parseval, a level change of $\pm2$ costs $(2\Qs)^2$ wherever it occurs. The energy depends on the quantizer only; it does not depend on the image, the prediction mode or the frequency.

\begin{lemma}[DC paths]
\label{lem:dc}
A flip in an Intra16$\times$16 LumaDC block or a ChromaDC block has, before rounding, the same total energy $4\Qs^2$ (with $\mathit{QP}_c$ for chroma), spread uniformly over the sixteen (resp.\ four) $4\times4$ blocks of the MB component.
\end{lemma}
\begin{proof}
LumaDC: the Hadamard output changes by $\pm2$ in all sixteen positions, and the DC scaling $(f\cdot16v_{m,0}2^{\lfloor\mathit{QP}/6\rfloor})/64$ turns this into a change of $v_{m,0}2^{\lfloor\mathit{QP}/6\rfloor}/2$ in the DC coefficient of each block, i.e.\ a constant residual change of $v_{m,0}2^{\lfloor\mathit{QP}/6\rfloor}/128$ over each block; the total energy is $256\,(v_{m,0}2^{\lfloor\mathit{QP}/6\rfloor}/128)^2=4\Qs^2$. ChromaDC: the $2\times2$ Hadamard output changes by $\pm2$ in four positions, the scaling $(f\cdot16v_{m,0}2^{\lfloor\mathit{QP}_c/6\rfloor})/32$ gives $v_{m,0}2^{\lfloor\mathit{QP}_c/6\rfloor}$ per block, and four blocks of sixteen samples give the same total.
\end{proof}

\textbf{Additivity.} The schedule uses one candidate per block, so two flips in the same block can only be a DC-path flip and an AC flip of the same MB component. Their errors are orthogonal (an AC basis function has zero mean, a DC change is constant over the block), so the direct energies add exactly. Rounding by $(\cdot+32)\gg6$ adds a bounded perturbation; in the demo block of Appendix~\ref{app:example} the exact rounded energy is 128 against a model value of 121.

\subsection{Propagation}
Ignoring clipping and the rounding of prediction, intra decoding is affine in the residual: in decoding order every block is its prediction (a linear function of earlier reconstructed samples) plus its residual. A direct error $\delta_k$ of flip $k$ therefore becomes a total error $\mu_k=(I-\mathcal P)^{-1}\delta_k$, where $\mathcal P$ is the strictly lower-triangular prediction operator of the frame, and deblocking adds a further, data-dependent smoothing. We define the \emph{propagation gain}
\begin{equation}
G_k=\frac{\|\mu_k\|^2}{4\Qs(\mathit{QP}_k)^2},
\end{equation}
which captures drift through intra prediction and deblocking ($G_k\approx\kappa_k$ when nothing propagates; rounding and clipping can push it slightly below). For a set of flips the total error is $\sum_k\mu_k$; cross terms $\langle\mu_k,\mu_l\rangle$ carry the signs of the flipped levels, which are those of the cover coefficients, and are zero-mean over the random keyed choice of carriers. Hence
\begin{equation}
\mathbb E[\SSE]\approx\sum_k 4\,\Qs(\mathit{QP}_k)^2\,G_k .
\end{equation}
$G_k$ depends on the prediction structure---MB type, prediction modes, block position, whether deblocking is on---which is why it is measured rather than derived (Section~\ref{sec:eval}).

\subsection{Forward model and inverse}
With constant QP, $N$ flips in the luma plane and a frame-level gain $\bar G$ (the geometric mean of $G_k$ fitted on data),
\begin{equation}
\PSNR_Y\approx10\log_{10}\frac{255^2\,W H}{4\,\bar G\,N\,\Qs(\mathit{QP})^2}.
\label{eq:forward}
\end{equation}
Three rules of thumb follow: doubling the number of flips costs $3.01$\,dB, raising QP by six costs $6.02$\,dB, and four times more pixels gain $6.02$\,dB at the same number of flips. Inverting Eq.~\eqref{eq:forward} for a target $T$:
\begin{equation}
N_{\max}(T)=\frac{255^2\,W H\,10^{-T/10}}{4\,\bar G\,\Qs(\mathit{QP})^2},\qquad
\mathit{cap}(T)=\frac{N_{\max}(T)}{p\,f_Y},
\label{eq:inverse-cap}
\end{equation}
where $p=1/2$ is the probability that a whitened bit changes the sign and $f_Y$ the share of luma candidates (a uniformly random schedule selects luma carriers at that rate). Chroma planes have their own budget with $\mathit{QP}_c$ and the chroma gain.

\textbf{Receiver computability.} A per-IDR cap is usable only if the receiver derives the same value from the stego stream. Eq.~\eqref{eq:inverse-cap} uses $W$, $H$, the slice QP, $f_Y$ and a protocol constant $\bar G$; by Proposition~\ref{prop:syntax} all of them are identical in the cover and the stego stream. An adaptive cap therefore needs no side information. The same constraint applies to any cost-aware schedule: costs may depend on positions, categories, MB types, prediction modes and coefficient magnitudes, but never on the sign values that carry data.

\subsection{Towards minimal-distortion embedding}
Eq.~\eqref{eq:flip} supplies an additive per-carrier cost $\rho_k=4\Qs(\mathit{QP}_k)^2\kappa_k\hat G_k$, with $\hat G_k$ predicted from invariant features (e.g.\ category and MB type). For a payload of $L$ bits, the distortion-limited sender of~\cite{filler2011stc,fridrich2007practical} changes carrier $k$ with probability $p_k=1/(1+e^{\lambda\rho_k})$, where $\lambda$ solves $\sum_k h(p_k)=L$ ($h$ the binary entropy) and the expected distortion is $\sum_kp_k\rho_k$; syndrome-trellis codes come close to this bound. Even without costs, matrix embedding reduces flips per bit: the present channel spends $1/2$ expected flip per bit, while a Hamming code with $r=3$ embeds three bits into seven carriers with $1-2^{-3}=0.875$ expected flips, i.e.\ $0.29$ flips per bit, a $10\log_{10}(0.5/0.29)\approx2.3$\,dB gain at equal payload. Candidates are plentiful (thousands per IDR against 64 used), so the larger carrier count is affordable. We leave syndrome coding to future work. A simple integer version of the cost, the drift tier of Eq.~\eqref{eq:tier}, is implemented as the \texttt{low-drift} selection policy (Section~\ref{sec:params}) and evaluated end to end in Section~\ref{sec:eval-compare}.
